%------------------------
% Resume Template
% Author : Anubhav Singh
% Github : https://github.com/xprilion
% License : MIT
%------------------------

\documentclass[a4paper,20pt]{article}

\usepackage{latexsym}
\usepackage[empty]{fullpage}
\usepackage{titlesec}
\usepackage{marvosym}
\usepackage[usenames,dvipsnames]{color}
\usepackage{verbatim}
\usepackage{enumitem}
\usepackage{hyperref}
\hypersetup{hidelinks}
\usepackage{fancyhdr}
\usepackage[UTF8, fontset=fandol]{ctex}

\pagestyle{fancy}
\fancyhf{} % clear all header and footer fields
\fancyfoot{}
\renewcommand{\headrulewidth}{0pt}
\renewcommand{\footrulewidth}{0pt}

% Adjust margins
\addtolength{\oddsidemargin}{-0.530in}
\addtolength{\evensidemargin}{-0.375in}
\addtolength{\textwidth}{1in}
\addtolength{\topmargin}{-.45in}
\addtolength{\textheight}{1in}

\urlstyle{rm}

\raggedbottom
\raggedright
\setlength{\tabcolsep}{0in}

% Sections formatting
\titleformat{\section}{
  \vspace{-10pt}\scshape\raggedright\large
}{}{0em}{}[\color{black}\titlerule \vspace{-6pt}]

%-------------------------
% Custom commands
\newcommand{\resumeItem}[2]{
  \item\small{
    \textbf{#1}{: #2 \vspace{-2pt}}
  }
}

\newcommand{\resumeItemWithoutTitle}[1]{
  \item\small{
    {\vspace{-2pt}}
  }
}

\newcommand{\resumeSubheading}[4]{
  \vspace{-1pt}\item
    \begin{tabular*}{0.97\textwidth}{l@{\extracolsep{\fill}}r}
      \textbf{#1} & #2 \\
      \textit{#3} & \textit{#4} \\
    \end{tabular*}\vspace{-5pt}
}

\newcommand{\resumeSubItem}[2]{\resumeItem{#1}{#2}\vspace{-3pt}}

\newcommand{\resumeProjectItem}[3]{
  \item
    \begin{tabular*}{0.97\textwidth}{l@{\extracolsep{\fill}}r}
      \textbf{#1} & \textit{#2} \\
    \end{tabular*}
    \vspace{-2pt}
      #3
    \par\vspace{-2pt}
}

\renewcommand{\labelitemii}{$\circ$}

\newcommand{\resumeSubHeadingListStart}{\begin{itemize}[leftmargin=*]}
\newcommand{\resumeSubHeadingListEnd}{\end{itemize}}
\newcommand{\resumeItemListStart}{\begin{itemize}}
\newcommand{\resumeItemListEnd}{\end{itemize}\vspace{-5pt}}

%-----------------------------
%%%%%%  CV STARTS HERE  %%%%%%

\begin{document}

%----------HEADING-----------------
\begin{tabular*}{\textwidth}{l@{\extracolsep{\fill}}r}
   \textbf{{\LARGE 聂家明}} \\
   \href{https://www.linkedin.com/in/jmnie/}{LinkedIn: in/jmnie}  & 邮箱:\href{mailto:jiaming.nie13@gmail.com}{jiaming.nie13@gmail.com} \\
  & 电话: (+86)-135-2171-3097
\end{tabular*}

\vspace{4pt}
\section{技术栈}
	\resumeSubHeadingListStart
    \resumeSubItem{后端开发}{~~~Python, FastAPI}
    \resumeSubItem{数据存储}{~~~MongoDB, Redis / Amazon MemoryDB}
    \resumeSubItem{基础设施}{~~~AWS, Docker, Apache Airflow, Jenkins CI/CD}
    \resumeSubItem{可观测性}{~~~Prometheus, Grafana}
	\resumeSubHeadingListEnd

\vspace{4pt}
\section{工作经历}
  \resumeSubHeadingListStart
  \vspace{4pt}
  \resumeSubheading{Youlify}{远程}
    {软件工程师（早期核心成员）}{2025.1 - 2026.9}
    \resumeItemListStart
        \item{开发医疗计费后端服务与 \textbf{Agent/LLM 工作流}，覆盖 \textbf{EOB} 解析、拒赔分析、申诉信生成、医疗编码审核和任务状态跟踪。}
        \item{为支持多诊所同时从 EHR 拉取患者资料和临床文档，将同步任务迁移至 \textbf{Airflow}，并通过 \textbf{Redis} 排队避免同一诊所的浏览器操作相互干扰；不同诊所可并行处理，用户可查看进度或取消任务。}
        \item{开发多个 \textbf{EHR 系统}的集成，使用 \textbf{Playwright} 与 API 同步患者、保险、就诊记录及文档，接入医疗计费流程。}
        \item{实现跨鉴权与业务服务的诊所级多角色权限、用户停用和会话撤销；在移除访问权限时支持理赔任务转交，并保留历史操作人记录。}
    \resumeItemListEnd

    \vspace{4pt}
  \resumeSubheading{北京峥研软件有限责任公司（peerup）}{上海, 中国}
    {软件开发工程师}{2023.6 - 2025.1}
    \resumeItemListStart
        \item{主导笔记 AI 助手的 \textbf{FastAPI} 后端建设，支持 \textbf{SSE} 流式响应、文件问答、复杂指令执行和编辑器内动作触发。}
        \item{设计基于 \textbf{Milvus} 与 \textbf{Elasticsearch} 的混合 \textbf{RAG}，整合公共知识库、用户内容和外部数据源，并通过 \textbf{Redis} 管理状态与用户级检索权限。}
        \item{维护 \textbf{Docker} 容器化与自动化 \textbf{CI/CD} 流水线，支撑 AI 功能稳定发布与持续迭代。}
      \resumeItemListEnd

    \vspace{4pt}
    \resumeSubheading{Oqton}{上海, 中国}
    {软件开发工程师}{2021.10 - 2023.2}
    \resumeItemListStart
        \item{面向工业制造场景开发 \textbf{Python} 服务，通过 API 完成机器数据采集与控制指令下发，参与企业级系统集成和现场交付。}
        \item{在 \textbf{Linux} 网关上部署和运维 \textbf{Docker} 应用，排查跨 \textbf{Elasticsearch} 与 \textbf{Kafka} 系统的生产环境问题。}
      \resumeItemListEnd

    \vspace{4pt}
    \resumeSubheading
		{微福思软件科技有限公司(Micro Focus)}{上海, 中国}
		{软件开发工程师}{2020.7 - 2021.9}
		\resumeItemListStart
              \item{使用 \textbf{JavaScript}、\textbf{Node.js} 与 \textbf{C++} 构建 Web 组件及跨域检测自动化能力，并优化内部案例处理工作流。}
		\resumeItemListEnd

\resumeSubHeadingListEnd

%-----------EDUCATION-----------------
\vspace{2pt}
\section{教育背景}
  \resumeSubHeadingListStart
    \resumeSubheading
      {伍斯特理工学院}{马萨诸塞，美国}
      {硕士 - 计算机科学}{2017.8 - 2019.12}
    \resumeSubheading
      {利物浦大学}{利物浦，英国}
      {学士 - 电气工程}{2015.8 - 2017.6}
    \resumeSubheading
      {西交利物浦大学}{苏州，中国}
      {学士 - 电气工程}{2013.8 - 2015.8}
    \resumeSubHeadingListEnd

%-----------PROJECTS-----------------
\newpage
\vspace{-5pt}
\section{技术项目}
\resumeSubHeadingListStart

\vspace{4pt}
\resumeProjectItem{医疗计费编码抽取与模型微调}{2026.2 - 2026.4}{
\resumeItemListStart
    \item{整理包含 \textbf{2.25 万条医疗编码记录}的微调训练数据集，使用固定随机种子划分训练集、验证集和测试集；最终使用 \textbf{3,900 条经上下文长度筛选的样本}进行混合数据监督微调。}
    \item{通过托管式微调接口编排\textbf{监督微调}任务，实现训练文件上传、任务提交、状态轮询、模型检查点跟踪与实验元数据管理自动化。}
    \item{建立分层离线评估与失败案例分析流程，分别评估\textbf{医疗操作编码（CPT）}、\textbf{诊断编码（ICD-10）}、修饰符及包含计费单位数的完整编码组合，计算精确率、召回率及 \textbf{Macro F1、Micro F1}。}
    \item{在同一独立留出的门诊测试集上，相较通用大模型基线，将 \textbf{CPT Macro F1 从 0.666 提升至 0.817}，将\textbf{完整编码组合严格匹配的 Macro F1 提升 0.218}；同时比较模型各检查点的效果与词元使用成本，辅助模型选择。}
\resumeItemListEnd
}

\vspace{4pt}
\resumeProjectItem{多 EHR 集成与临床数据同步}{2025.1 - 2026.9}{
\resumeItemListStart
    \item{集成 Athena、MDLand、Practice Fusion 等多个 EHR，使用 \textbf{Playwright} 和接口采集患者资料、保险、就诊记录及临床文档。}
    \item{实现自动登录、多因素认证、每日同步与历史数据补录，并在支持接口的 EHR 中用接口调用替代浏览器操作。}
    \item{通过独立的\textbf{租户数据微服务}接口，按诊所保存患者、保险及收费记录，将采集数据接入医疗计费流程。}
    \item{将临床文档关联至对应患者和就诊记录，供理赔审核使用；增加数据校验与错误记录，便于排查和补齐缺失数据。}
\resumeItemListEnd
}

\vspace{4pt}
\resumeProjectItem{基于 Airflow 的数据同步与任务管理平台}{2026.6 - 2026.9}{
\resumeItemListStart
    \item{为统一调度患者资料与临床文档同步，将集成服务中的后台任务迁移至 \textbf{Airflow} 工作流，保持原有业务接口和任务状态定义不变。}
    \item{使用 \textbf{Redis/MemoryDB 有序集合}实现按环境、诊所隔离的先进先出队列，使同一诊所的浏览器任务依次执行、不同诊所并行处理；排队期间不占用执行资源，直接调用接口的任务无需等待浏览器队列。}
    \item{通过定期心跳确认任务仍在运行，并在任务超时失效或被取消时清理队列，避免执行进程中断后留下的任务持续阻塞后续同步。}
    \item{在任务分发前将记录保存至 \textbf{MongoDB}，并在工作流结束时写入最终状态；写入失败时仅重试状态保存，避免重复执行浏览器操作，替代原先每分钟核对任务状态的方式。}
\resumeItemListEnd
}

\vspace{4pt}
\resumeProjectItem{后端基础设施与运维}{2025.5 - 2026.9}{
\resumeItemListStart
    \item{主导单体后端按业务领域拆分为 \textbf{FastAPI} 服务，覆盖 Portal API、诊所管理、鉴权和 GenAI，使各服务能够独立部署。}
    \item{为多个后端仓库建立统一的 \textbf{Docker/Jenkins 持续集成与部署流程}，串联合并请求检查、镜像构建和多环境发布；将镜像版本关联至代码提交，支持各服务独立发布与版本追踪。}
    \item{开发 \textbf{Google Chat 部署机器人}，支持 Jenkins 环境选择、确认及各环境部署状态跟踪；持久化部署状态，防止重复触发并在重启后恢复跟踪。}
    \item{接入 \textbf{Prometheus/Grafana} 服务监控，跟踪请求错误与延迟；排查容器启动、日志依赖及 Nginx 路由故障，恢复开发与测试环境服务。}
    \item{通过共享客户端、连接池上限和超时配置完善 \textbf{MongoDB} 连接管理；将只读报表与定时任务查询分流至 secondary 副本，减轻主节点读负载。}
\resumeItemListEnd
}

\resumeSubHeadingListEnd

\end{document}
